% Source: https://artofproblemsolving.com/community/c3h2258725_sierpinski_triangle
% Stable identity: aops-topic-2258725-q01
Let $T_n$ be the number of triangles in level $n$; $T_1=1$ and $T_2=5$.

\textbf{Structure.} Let the big triangle be $ABC$ and let $D,E,F$ be the midpoints of $BC,CA,AB$. Level $n$ ($n\ge2$) consists of three copies of level $n-1$, scaled by $\dfrac{1}{2}$, in the corner triangles $AFE$, $FBD$ and $EDC$. These copies meet only at the points $D$, $E$, $F$, and the drawn segments of each copy lie in its corner triangle. Every drawn segment is parallel to a side of $ABC$, so every triangle in the figure is equilateral and points either up (like $ABC$) or down.

For example, level $3$ consists of three copies of level $2$ (shaded) and the red middle triangle $DEF$:

\begin{asy}
usepackage("amsmath");
size(5cm);
defaultpen(fontsize(10pt));
void sier(pair a, pair b, pair c, int d, pen p){
  draw(a--b--c--cycle,p);
  if(d>1){ pair ab=(a+b)/2, bc=(b+c)/2, ca=(c+a)/2;
    sier(a,ab,ca,d-1,p); sier(ab,b,bc,d-1,p); sier(ca,bc,c,d-1,p); }
}
pair B=(0,0), C=(4,0), A=(2,2*sqrt(3)), D=(B+C)/2, E=(C+A)/2, F=(A+B)/2;
fill(A--F--E--cycle,paleyellow); fill(F--B--D--cycle,palecyan); fill(E--D--C--cycle,palegreen);
sier(A,F,E,2,black); sier(F,B,D,2,black); sier(E,D,C,2,black);
draw(A--B--C--cycle,linewidth(1));
draw(D--E--F--cycle,red+linewidth(1.2));
label("$A$",A,N); label("$B$",B,SW); label("$C$",C,SE);
label("$D$",D,S,red); label("$E$",E,NE,red); label("$F$",F,NW,red);
\end{asy}

\textbf{Claim.} Every triangle in level $n$ other than $ABC$ and the middle triangle $DEF$ lies inside one of the three copies.

Let $T$ be a triangle in the figure that is not inside one copy. Walking around the boundary of $T$, we can pass from one copy to another only through $D$, $E$ or $F$, and the walk passes each point once. If the boundary used segments of only two copies, say those at $A$ and $B$, it would have to go from one to the other and back through $F$, the only common point of these copies, twice, which is impossible. So the boundary of $T$ uses segments of all three copies and passes through $D$, $E$ and $F$. Also $T$ lies inside $ABC$, since all drawn segments do.
\begin{itemize}
\item If $T$ points up, it has a bottom side parallel to $BC$ and two sides parallel to $AB$ and $AC$. As $T$ lies inside $ABC$ and contains the point $D$ of line $BC$, its bottom side lies on line $BC$. As $T$ lies inside $ABC$ and contains the point $E$ of line $AC$, its side parallel to $AC$ lies on line $AC$; in the same way (using $F$) its side parallel to $AB$ lies on line $AB$. So $T=ABC$.
\item If $T$ points down, the only point of $T$ at its lowest level is its bottom vertex. As $T$ lies inside $ABC$ and contains the point $D$ of line $BC$ (the lowest level of $ABC$), its bottom vertex is $D$. The two sides through $D$ are parallel to $AB$ and $AC$, so they lie on the midlines $DE$ and $DF$. These lines leave triangle $ABC$ at $E$ and $F$, so $T$ reaches no higher than the level of $EF$; as $T$ contains $E$, its top side lies on line $EF$. So $T=DEF$.
\end{itemize}
This proves the claim.

\textbf{Recursion.} Each copy contains $T_{n-1}$ triangles, the copies share no triangle, and we add $ABC$ and $DEF$:
\[T_n=3T_{n-1}+2.\]
Adding $1$ to both sides, $T_n+1=3(T_{n-1}+1)$, so $T_n+1=3^{n-1}(T_1+1)=2\cdot3^{n-1}$ and
\[T_n=\boxed{2\cdot3^{n-1}-1}.\]
This gives $1,5,17,53,161,\dots$ for $n=1,2,3,4,5,\dots$
